\documentclass{article}
\usepackage{geometry}
\usepackage{graphicx}
\usepackage{filecontents}

% --- 模拟外部子文件 ---
\begin{filecontents}{sub_chapter.tex}
\section{Complex Layout Analysis}
Here we utilize a sub-file to test the file-tracking capability of LPSB.

\begin{figure}[h]
    \centering
    % 使用 resizebox 强行缩放，这是 tagpdf 的噩梦，但 LPSB 能搞定
    \resizebox{0.5\textwidth}{!}{%
        \includegraphics[width=5cm]{example-image-a}%
    }
    \caption{A Resized Image in a Subfile}
    \label{fig:resized_sub}
\end{figure}
\end{filecontents}
% ---------------------

\begin{document}

\section{Introduction}
Welcome to the LPSB Demo. This demonstrates the LaTeX-PDF Semantic Bridge system.
\label{sec:intro}

A second paragraph here shows how paragraph tracking works.

\begin{figure}[t]
    \centering
    \includegraphics[width=3cm]{example-image-b}
    \caption{Standard Figure}
    \label{fig:std_fig}
\end{figure}

As shown in Figure~\ref{fig:std_fig}, this is how cross-references work.

% 引入子文件
\input{sub_chapter}

\section{Table Analysis}
\begin{table}[h]
    \centering
    \caption{Data Summary}
    \begin{tabular}{|c|c|}
        \hline
        ID & Value \\
        \hline
        1 & 99.9 \\
        \hline
    \end{tabular}
    \label{tab:data}
\end{table}

Table~\ref{tab:data} shows our results. See also Figure~\ref{fig:resized_sub}.

\subsection{Detailed Results}
This subsection demonstrates the section hierarchy tracking.

\subsubsection{Even More Detail}
Nested sections are properly handled by the section tree builder.

\section{Conclusion}
Final thoughts and summary.

\end{document}
